\chapter*{چکیده}

تست خودکار رابط کاربری برای نرم‌افزارهای مبتنی بر
\lr{Qt C++}
تا امروز عمدتاً به ابزارهای تجاری محدود بوده و ابزارهای متن‌باز موجود، تنها زیرساخت خام (درایور و پروتکل) را فراهم می‌کنند. این پایان‌نامه طراحی و پیاده‌سازی یک چارچوب تست سرتاسری
\lr{\gls{E2E}}
متن‌باز برای برنامه‌های دسکتاپ مبتنی بر
\lr{Qt}
در بستر ویندوز را گزارش می‌کند. چارچوب بر زیرساخت
\lr{WinAppDriver}/\lr{Appium}
استوار است و با پنج لایه‌ی معماری \pdash{} مکان‌یابی با نهانگاه دسته‌ها و جستجوی محدوده‌ای، اشیای صفحه، زبان مخصوص دامنه، کارخانه‌ی درایور با اتصال فازی، و جداسازی داده و گزارش‌گیری \pdash{} پایداری و خوانایی تست‌ها را فراهم می‌کند. دو قرارداد پایداری (شناسه پایدار و قرارداد خوانش) رابطه‌ی برنامه‌ی هدف و چارچوب را قواعد‌مند کرده و تست نمودارهای گرافیکی را به تست داده‌ی خوانا تبدیل می‌کند.

برای ارزیابی، برنامه‌ی نمونه‌ی «ردیاب سیگنال ایران» (مبتنی بر
\lr{Qt}
و
\lr{SQLite})
ساخته شد که جریان‌های دشوار تست شامل ورود چندپنجره‌ای، فیلترهای آبشاری، صفحه‌بندی، مجازی‌سازی جدول، نمودار و گزارش‌گیری را داراست. ۲۵ سناریوی سرتاسری روی آن پیاده شد که اجرای کاملشان پایدار، بازتولیدپذیر و در حدود شش دقیقه است و گزارش تعاملی
\lr{Allure}
با پیوست تصویر خطا تولید می‌شود. مسیر انتقال چارچوب به پایپلاین‌ی
\lr{\gls{CI}}/\lr{\gls{CD}}
نیز طراحی شده است.

\textbf{کلیدواژه‌ها:} تست سرتاسری، خودکارسازی رابط کاربری، \lr{Qt}، \lr{WinAppDriver}، \lr{Appium}، زبان مخصوص دامنه، شیء صفحه